\documentclass[11pt]{article}
\usepackage[margin=1in]{geometry}
\usepackage[T1]{fontenc}
\usepackage{lmodern}
\usepackage{amsmath,amssymb,booktabs,array}
\usepackage[numbers,sort&compress]{natbib}
\usepackage[colorlinks=true,allcolors=blue]{hyperref}
\usepackage{microtype}
\newcommand{\R}{\mathbb{R}}
\newcommand{\status}{Private exploratory working draft}
% Generated from immutable receipts by tools/summarize_pilots.py; do not hand-edit.
\newcommand{\HenonTrials}{65}
\newcommand{\HenonCycles}{39}
\newcommand{\TrainingTrials}{162}
\newcommand{\TrainingCycles}{0}
\newcommand{\TrainingUnresolved}{107}
\newcommand{\TrainingDivergent}{55}
\newcommand{\UtilityTrials}{96}
\newcommand{\PolicyTrials}{24}
\newcommand{\PolicyDivergent}{9}
\newcommand{\RecurrenceTriggers}{0}

\title{Do Periodic Windows Help Control Neural Networks?\\
\large An Exploratory Pilot and Protocol for Testing Shrimp Geometry}
\author{Authorship to be finalized by the human researcher}
\date{8 September 2026}

\begin{document}
\maketitle
\noindent\textbf{\status.} This manuscript reports initial local pilots and specifies follow-up questions and methods. It reports no confirmed neural-network shrimp and no demonstrated optimization or reasoning benefit. All pilot artifacts are exploratory. This document is not a preregistration or publication-ready article.

\medskip
\noindent\textbf{AI assistance and editorial status.} Generative-AI tools assisted with the research design, code, execution, local audits, source checking, and draft. The human selected the research question. Human shaping, artifact inspection, authorship, and final claim approval remain pending. The canonical publication disclosure will be applied at release only when its statements are true. No paid Pro review was requested or run.

\begin{abstract}
Shrimp-shaped periodic windows organize stable dynamics in some nonlinear systems, including small recurrent neural networks. Their presence alone does not imply useful computation. We separate existence of neural periodic orbits, verification of their parameter-space organization, and operational value. In initial local pilots, the numerical instrument identified attracting nontrivial cycles in \HenonCycles{} of \HenonTrials{} H\'enon grid cells. A tiny tanh training scan found \TrainingCycles{} such cycles among \TrainingTrials{} cells, with \TrainingUnresolved{} unresolved and \TrainingDivergent{} operationally divergent cases. The recurrence-triggered learning-rate policy never intervened in its \PolicyTrials{} runs and reproduced the fixed-rate trajectories exactly; all four policies had \PolicyDivergent{} divergent cases each. These results demonstrate neither neural shrimps nor a control advantage. They motivate a staged continuation: first identify and verify a neural periodic regime, then test whether its geometry improves intervention choices beyond equally funded loss, recurrence, and local-stability diagnostics on held-out tasks. Any claimed benefit must include mapping costs and retain failed cases. Latent reasoning remains a proposed extension.
\end{abstract}

\section{Why a periodic window might matter}
Repeated dynamics can be stable without being useful. A training run that reaches a nontrivial attracting cycle ceases to make asymptotic progress under its fixed update, yet its predictions may already be adequate. A recurrent solver can return the correct answer while its internal state keeps moving. Conversely, a stable but incorrect answer can waste an inference budget. The question is therefore operational: does knowledge of periodic-window organization improve quality, reliability, or resource use?

A \emph{shrimp} denotes a characteristic organization of periodic windows in a control-parameter plane, rather than a region of the evolving state space. Gallas described such domains for the H\'enon map \citep{gallas1993}. Neural-network existence is already established in a specific precedent: da Silva and Rech report shrimp structures in a three-neuron continuous-time Hopfield-type system with tanh activations \citep{dasilva2022}. That precedent does not identify their utility for learning.

Two recent motivations concern different spaces. Sohl-Dickstein's preprint and companion article examine training success across hyperparameters \citep{sohldickstein2024,sohldicksteinblog2024}. Lai and colleagues' preprint reports transient dynamics from slices of initial latent states in recurrent reasoning models \citep{lai2026}. Neither map, as presented, identifies the primitive periodic orbits and parameter bifurcations required by this study.

\paragraph{Prior work, proposed contribution, and non-claims.}
Oscillating optimization and a period-doubling route to chaos already appear in neural training research \citep{cohen2021,kalra2025}. Our proposed contribution is an incremental-value test of parameter-window geometry. We do not claim first neural shrimps, universal periodicity, a general need to suppress chaos, or any large-model speedup. No result in this draft supports those claims.

\section{State, controls, and utility are separate objects}
For a deterministic training problem $\mathcal{D}$, write
\begin{equation}
 \theta_{t+1}=F_\eta(\theta_t)
 =\theta_t-D_\eta\nabla L_{\mathcal D}(\theta_t),\qquad
 D_\eta=\operatorname{diag}(\eta_1 I,\eta_2 I).
\end{equation}
Weights $\theta$ are state variables; positive layer learning rates $\eta=(\eta_1,\eta_2)$ are controls. For latent inference, a fixed model and problem instead define $h_{t+1}=G_\alpha(h_t;x)$. The hidden state evolves while weights are controls. Every evolving cache, auxiliary latent, momentum variable, or schedule phase must be included when defining an autonomous map.

Continuous gradient flow has a restrictive property:
\begin{equation}
 \dot\theta=-\nabla L(\theta)
 \quad\Longrightarrow\quad
 \frac{dL}{dt}=-\|\nabla L(\theta)\|^2.
\end{equation}
It cannot support a nonconstant periodic orbit, since integrating the right-hand side over a period forces the gradient to vanish throughout. Finite-step gradient descent need not retain monotonic loss. This distinction permits a periodic-orbit search but does not guarantee shrimp organization.

The utility of a dynamical regime must be measured independently of its label. For training, candidate endpoints are validation error at a total computation budget and computation to a fixed quality threshold. For reasoning, they are correctness at a budget, incorrect termination, and latency. Attraction and periodicity are predictors, not utility definitions.

\section{From recurrence to verified parameter structure}
A periodic point of primitive period $p$ satisfies $F_\eta^p(\theta)=\theta$ without satisfying a smaller-period equation. Scalar loss recurrence is insufficient because many weight states can share a loss value. The pilot searches full-state tails over repeated cycles, checks a nonzero cycle amplitude relative to numerical tolerance, rejects smaller periods, and evaluates closure independently.

For a verified orbit $\theta_0,\ldots,\theta_{p-1}$, define its period Jacobian product
\begin{equation}
 M_p=DF_\eta(\theta_{p-1})\cdots DF_\eta(\theta_0).
\end{equation}
A spectral radius $\rho(M_p)<1$, separated from the threshold by numerical error, establishes local attraction of an isolated orbit. It does not measure basin size. Flat or symmetry-related directions can produce neutral multipliers; near-unit values require an explicit nonhyperbolic or unresolved status. Discarding them to manufacture stability is impermissible.

Verified periodic pixels are still insufficient to name a shrimp. A later stage must continue periodic branches and locate their saddle-node and period-doubling boundaries around connected isoperiodic regions. Longer horizons, changed precision, and multiple initial states assess transient misclassification and multistability. The initial H\'enon control tests the numerical instrument, not the validity of the neural conclusion.

The chosen parameter coordinates matter for geometric predictors. A nonlinear axis transformation changes widths, normals, and visual shape. The study must therefore freeze meaningful coordinates, such as log learning rates, and distinguish coordinate-dependent distances from bifurcation organization.

\section{Experiments and their permitted inferences}
\subsection{Initial exploratory implementation}
The initial target is deterministic full-batch training of a tiny tanh network over two layer learning rates. The network has two inputs, three hidden units, one scalar output, and biases in both layers. All thirteen parameters evolve. The objective is half mean squared error plus $0.001\|\theta\|^2/2$, including biases in the penalty. Predictive MSE excludes that penalty. Data come from a random tanh teacher with independent standard-normal input splits containing 24 training, 48 validation, and 96 test examples. Teacher metadata, all splits, and initial states are prepared before the freeze, bound by manifest hashes, and verified when loaded. This is inspired by the fractal-training construction rather than an exact replication.

The result-free manifests specify a 65-cell H\'enon control with 3,000 updates per cell and a $9\times9$ learning-rate training scan for seeds 101 and 102 with 2,000 updates per cell. Each neural seed defines both a task and its one initialization, so task and initialization variation are not separately estimated. Scans retain up to 512 tail states and screen periods up to 16 in float64. All requested cases, including divergence, unresolved states, and failed orbit checks, must remain in the census. Bounded aperiodic-looking tails are not automatically classified as chaos.

The first utility comparison uses four task/initialization seeds, 201 through 204, and six initial rate pairs, for 96 runs across four policies. The policies are fixed-rate training, predetermined scheduled decay, scalar stagnation-triggered decay, and recurrence-triggered decay. Each run has a maximum of 1,200 gradient updates. Every decay halves both rates, with at most two decays. The schedule acts at updates 400 and 800; adaptive rules inspect every 40 updates from update 200. Recurrence uses full-weight screening rather than online orbit certification. Each policy is executed separately, but paths are identical until the first differing intervention. Every generated state, pre-update regularized loss, realized learning rate, and trigger is retained.

The pilot's primary readouts are terminal test MSE and divergence fraction. Test and validation outcomes do not trigger interventions. Its descriptive report presents a policy mean only when all corresponding runs are finite, together with failures, triggers, costs, and seed-level records. It provides no inferential p-values or confidence intervals. This differs from the proposed fixed-resource validation endpoint of the later held-out study.

Both adaptive detectors wait at least two check intervals after a rate change. The recurrence detector uses only states from the current constant-rate segment, and stagnation uses two loss windows from that segment. The finite-history trigger is not an independent certificate of a periodic orbit.

This comparison can test whether recurrence is a useful diagnostic. It contains no intervention that uses shrimp geometry, so it cannot support a geometry-specific benefit. Equal maximum update budgets also do not establish equal total cost or a speedup; divergent runs stop early. Reported wall time includes common pilot instrumentation, not just an optimized policy implementation. Diagnostic overhead and memory costs require separate accounting. Pilot execution follows the full code and plan commit and push; no target outcomes are required to validate this design.

\subsection{Proposed held-out geometry experiment}
The main hypothesis is that geometry provides predictive and control information beyond simpler observables. This stage is not yet implemented or prospectively frozen. Its machine plan must bind the exact sample, features, thresholds, interventions, budgets, and analysis before held-out outcomes exist.

Use separate problems and seeds for discovery, controller development, and evaluation. A synthetic teacher or data-generating problem is the independent task unit. Nearby map pixels and repeated time points are not independent replications. Hold out whole windows or spatial blocks in addition to tasks to avoid reporting local interpolation as transfer.

The predictive comparison uses identical fitted models and tuning budgets. Baseline features include recent loss, progress, gradient and update norms, recurrence residual, and a finite-time stability diagnostic. The augmented model adds measured period, boundary distance, local width, and boundary direction. Its label is future task utility after a specified intervention, not recurrence itself.

The control comparison branches common checkpoints into a frozen action set. Required baselines include a predetermined schedule, stagnation, recurrence, local stability, and ordinary locally sampled performance maps with matched scan costs. Compare random intervention directions of equal norm with geometry-informed directions. Give comparators the same doses, opportunities, and tuning budgets. This distinguishes geometry from the effects of extra trials or a stronger action set.

\subsection{Costs, uncertainty, and failure retention}
The proposed primary endpoint is paired validation error at a fixed total resource budget. Its exact aggregation, divergence treatment, and practical-effect threshold must be frozen after the pilot and before held-out evaluation. A complementary target-attainment endpoint treats unsolved cases as budget-censored. Reasoning experiments instead require exact task correctness at equal total inference cost.

Charge map construction, controller tuning, failed probes, Jacobian calculations, state storage, transfers, and extra validation calls. Report cold-start cost and explicitly amortized cost for a declared number of future runs. The experimentally reused prefix of several policies is charged to each hypothetical deployed policy. No transfer of an atlas across tasks is assumed free.

Use paired task-level estimates with seed variation retained inside tasks. Freeze the cluster-bootstrap procedure, number of draws, primary contrast, multiplicity rule, and missingness handling before evaluation. Report intervention effects alongside the unselected prevalence of problematic cycles. Rescue rates in an enriched collection of traps do not establish population-level savings.

\section{A mechanistic extension to latent reasoning}
Stable repetition and transient delay require different interventions. The Lai et al. preprint uses decoded settling and reports transient structure; its appendices also describe excluded capped or inconsistent slices and distinguish slim fractals \citep{lai2026}. Our utility extension would retain failures and measure both decoded correctness and full-state recurrence.

A clean starting intervention is positive relaxation,
\begin{equation}
 G_\alpha(h)=(1-\alpha)h+\alpha F(h),\qquad \alpha>0.
\end{equation}
It preserves the fixed-point equation: $G_\alpha(h)=h$ if and only if $F(h)=h$. At a fixed point, each eigenvalue $\lambda$ of $DF$ becomes
\begin{equation}
 \mu=1-\alpha+\alpha\lambda.
\end{equation}
Hence positive relaxation cannot stabilize a real $\lambda>1$ direction, though it can change oscillatory stability. This is a local algebraic fact, not a guarantee of global accuracy or faster convergence. Relaxation can change basins, transients, and nontrivial cycles.

Two controls preserving the fixed-point equation could be relaxation and momentum, with the previous iterate included in the state. Feedback gain is another experimental control but generally changes the fixed-point equation, so it changes the solver and requires a different interpretation. Compare damping, restarts, and geometry-informed actions against residual and stagnation policies on untouched tasks, with matched total inference budgets.

\section{Initial pilot results}
The initial runs used the complete code and input revision
\href{https://github.com/tdj28/neural-shrimp/commit/bd3d733b51e9f1ae9247ba712633aca023e7afac}{\texttt{bd3d733b51e9}} after it was pushed to the private repository. All planned rows were retained. Local receipt verification checked file hashes and the exact trial census. Empirical counts in this manuscript are generated from evidence through \texttt{generated-pilot.tex}; the immutable raw attempts and descriptive reports are retained under \texttt{evidence/}. These are exploratory outcomes.

\begin{center}
\begin{tabular}{p{0.24\linewidth}p{0.65\linewidth}}
\toprule
Pilot & Observed outcome \\
\midrule
H\'enon control & \HenonCycles{} of \HenonTrials{} sampled cells had a numerically verified attracting nontrivial periodic orbit. These are cell counts, not counts of distinct windows. \\
Neural training scan & \TrainingCycles{} verified attracting nontrivial cycles among \TrainingTrials{} cells; \TrainingUnresolved{} unresolved and \TrainingDivergent{} operationally divergent. \\
Policy comparison & \UtilityTrials{} runs across four policies; each policy had \PolicyDivergent{} divergent cases out of \PolicyTrials{}. The recurrence policy triggered in \RecurrenceTriggers{} runs. \\
\bottomrule
\end{tabular}
\end{center}

The classical control shows that the instrument can find and numerically verify periodic orbits in the tested H\'enon region. It does not certify shrimp topology or validate an unobserved neural orbit. The neural scan is unresolved with respect to existence outside the tested settings: a finite grid and horizon cannot establish global absence, and unresolved tails are not certified chaos.

For every one of the \PolicyTrials{} task--rate pairs, an independent local comparison found that fixed-rate and recurrence-policy files contained exactly equal saved state, loss, and learning-rate arrays, with NaNs compared at matching positions. Their statuses, completed update counts, and terminal predictive MSE values also matched. The recurrence detector therefore added observation overhead without changing the dynamics in this sample. This is evidence of a non-intervening pilot, not an equivalence test of effective intervention methods.

Every policy retained divergent cases, so the prespecified descriptive report withholds a successful-only mean test MSE. No inferential significance or net-compute advantage is claimed. The runs used local CPU resources with no paid cloud computation; recorded paid compute and paid review spend for these pilots was zero dollars. Local wall time is not a total economic-cost estimate.

\paragraph{What these outcomes change.}
The next scientific step is regime identification before benefit testing. A separately documented local successor should identify a reproducible nontrivial neural periodic regime, use longer horizons or targeted parameter refinement to distinguish drift and transients, and verify any candidate on the complete state. A published small-model precedent can serve as a second positive control without being described as a new discovery. Only after such a regime is verified should a controlled utility experiment measure when leaving, entering, or remaining in it helps task performance. Further search and controller development remain exploratory; neither the current empty scan nor the inactive trigger warrants paid scale-up.

\section{Discussion}
This section develops hypotheses beyond the current evidence. The initial pilots did not establish a neural shrimp or a useful recurrence intervention: the neural scan found no verified attracting nontrivial cycle, and the recurrence policy never acted. The plausible practical use remains a decision aid: recognize a costly dynamical regime and choose an intervention with a favorable quality--cost tradeoff. The distinctive contribution, if one exists, would come from parameter organization predicting which direction to move and how robust the destination is.

There are several informative future failure outcomes. If ordinary recurrence matches an active geometry method, the simpler diagnostic may be sufficient for the tested population. If mapping overhead erases per-run gains, there is no demonstrated net benefit at that deployment scale. If an intervention accelerates training loss reduction while harming validation quality, both effects must be reported. The current finite search constrains only its specified range, horizons, precision, and initialization choices.

There is no blanket prescription to seek or avoid these regions. For a task requiring a single fixed answer, an incorrect attracting cycle is a candidate failure mode. A stable cycle with correct readout could instead be harmless or support a stopping rule. Transient sensitivity may be productive. Those possibilities are separated by measured task outcomes, not by the appearance of a parameter-space picture.

\section*{Reproducibility and research status}
The private repository contains executable pilots, frozen inputs and configurations, immutable raw attempts, descriptive reports, a research plan, a claim ledger, and reference provenance. The initial program is exploratory and has no registry record. Private Git history supports internal chronology but is not public prospective verification. Any eventual confirmatory study requires a fresh frozen design, preserved failures and amendments, complete cost receipts, and outcome-linked claims. Human editorial review and release decisions remain pending; routine local audits do not require paid Pro review.

\bibliographystyle{plainnat}
\bibliography{references}
\end{document}
